\documentclass[11pt]{article}
\usepackage{cs152}

%%%%%%%%%%%%%%%%%%%% name/id
\rfoot{\small Brian Park | 3033045069}


%%%%%%%%%%%%%%%%%%%% Course/HW info
\newcommand*{\instr}{John Wawrzynek}
\newcommand*{\term}{Spring 2022}
\newcommand*{\coursenum}{CS 152}
\newcommand*{\coursename}{Computer Architecture and Engineering}
\newcommand*{\hwnum}{2}

%%%%%%%%%%%%%%%%%%%%%%%%%%%%%% Document Start %%%%%%%%%%%%%%%%%
\begin{document}

%%%%%%%%%%%%%%%%%%%%%%%%%%%%%%%%%%%%%%%%%%%%%%%%%%%%%%%%%%%%%%%%%%%%%%%%%%%%%%%%%%%%%%%%
% Question 3
%%%%%%%%%%%%%%%%%%%%%%%%%%%%%%%%%%%%%%%%%%%%%%%%%%%%%%%%%%%%%%%%%%%%%%%%%%%%%%%%%%%%%%%%
\section*{(Directed) Problem 3: Matrix Transposition Case Study}


%%%%%%%%%%%%%%%%%%%%%%%%%%%%%%%%%%%%%%%%%%%%%%%%%%%%%%%%%%%%%%%%%%%%%%%%%%%%%%%%%%%%%%%%
% Question 3.3
%%%%%%%%%%%%%%%%%%%%%%%%%%%%%%%%%%%%%%%%%%%%%%%%%%%%%%%%%%%%%%%%%%%%%%%%%%%%%%%%%%%%%%%%
\subsection*{3.3 Matrix Transposition Case Study}
\begin{Parts}
	\Part How many cycles does the transpose operation take? \\
		The transpose operation takes 1350476 cycles.
	\Part What is the miss rate of the L1 data cache? \\
		The miss rate of the L1 data cache is $\frac{18661}{16384 + 16407}$. The total number of memory accesses is found by counting the number of loads and stores. Therefore, the miss rate is 56.9\%.
	\Part Why does the naive transpose code result in non-ideal cache performance? \\
		The code has to transfer data diagonally through the matrix. If the matrix being transposed is fairly large, then a large stride access can hurt cache hits and memory performance.
	\Part Which memory access in the code incurs the most misses, and why? \\
		Obviously, this line of code: \verb|dst[j][i] = src[i][j];| incurs the most misses. Specifically, writing to \verb|dst[j][i]| will incur the most since the CPU will have already loaded \verb|src[i][j]| to a register from memory, but when it wants to write to \verb|dst[j][i]|, it will most likely miss the cache and stall for a memory fetch. 
\end{Parts}
\newpage

%%%%%%%%%%%%%%%%%%%%%%%%%%%%%%%%%%%%%%%%%%%%%%%%%%%%%%%%%%%%%%%%%%%%%%%%%%%%%%%%%%%%%%%%
% Question 3.4
%%%%%%%%%%%%%%%%%%%%%%%%%%%%%%%%%%%%%%%%%%%%%%%%%%%%%%%%%%%%%%%%%%%%%%%%%%%%%%%%%%%%%%%%
\subsection*{3.4 Cache Blocking}
\begin{Parts}
	\Part For the given matrix dimensions, what is the optimal blocking factor B? \\
	The optimal blocking is $8\times8$. Also, some loop ordering needed to be swapped in order for this to be cache optimal.
	\Part What is the performance improvement using cache blocking over the naive code? \\
	This achieves a miss rate of 21.21\%, which is an improvement over naive code by 2.68 times.
	\Part It turns out that the block size B×B which yields the lowest miss rate is much smaller than what one might expect based solely on the 4 KiB capacity of the L1 data cache. What is the reason for this? (Hint: Consider the access pattern within a block, particularly how the rows of the rectangular matrices map to cache sets and which type of cache misses dominate for larger B.) \\
	Considering that the L1 cache is 4KiB, this can hold 1024 32 bit integers. Since the matrix to be transposed will be accessed in column major, this will incur much more cache misses. That is why the block size is much smaller.
\end{Parts}
\newpage

%%%%%%%%%%%%%%%%%%%%%%%%%%%%%%%%%%%%%%%%%%%%%%%%%%%%%%%%%%%%%%%%%%%%%%%%%%%%%%%%%%%%%%%%
% Question 3.5
%%%%%%%%%%%%%%%%%%%%%%%%%%%%%%%%%%%%%%%%%%%%%%%%%%%%%%%%%%%%%%%%%%%%%%%%%%%%%%%%%%%%%%%%
\subsection*{3.5 Cache Parameters}
\begin{Parts}
	\Part How does performance and miss rate change when associativity is increased? \\
	When the associativity is increased, the miss rate worsens. The table below shows the relationship between the number of ways in the cache and the miss rate.

\begin{table}[ht]
\caption{$N$-Way Set Associativity vs. Miss Rate} % title of Table
\centering % used for centering table
\begin{tabular}{c c} % centered columns (4 columns)
\hline\hline %inserts double horizontal lines
$N$ & Miss Rate \\ [0.5ex] % inserts table
%heading
\hline % inserts single horizontal line
1-way & 21.21\% \\
2-way & 24.54\% \\
4-way & 28.05\% \\
8-way & 31.72\% \\
[1ex] % [1ex] adds vertical space
\hline %inserts single line
\end{tabular}
\label{table:nonlin} % is used to refer this table in the text
\end{table}
	

	\Part Explain why higher associativity is or is not beneficial for this particular kernel. \\
	A higher associativity is not beneficial since transposing takes advantage of spatial locality.
\end{Parts}
\newpage

%%%%%%%%%%%%%%%%%%%%%%%%%%%%%%%%%%%%%%%%%%%%%%%%%%%%%%%%%%%%%%%%%%%%%%%%%%%%%%%%%%%%%%%%
% Question 3.6
%%%%%%%%%%%%%%%%%%%%%%%%%%%%%%%%%%%%%%%%%%%%%%%%%%%%%%%%%%%%%%%%%%%%%%%%%%%%%%%%%%%%%%%%
\subsection*{3.6 Multi-level Caches}
\begin{Parts}
	\Part Does adding another level of cache blocking improve performance compared to your previous code from 3.5? Why or why not? (Hint: Consider whether there is any locality left for the L2 cache to exploit.) \\
\end{Parts}

\newpage
\section*{(Open-ended) Problem 4.1: Validation and Reverse Engineering of Memory Hierarchies}
\section*{(Open-ended) Problem 4.2: Design Your Own Hardware Prefetcher}
\section*{(Open-ended) Problem 4.3: Design Your Own Replacement Policy and Victim Cache}
\newpage

%%%%%%%%%%%%%%%%%%%%%%%%%%%%%%%%%%%%%%%%%%%%%%%%%%%%%%%%%%%%%%%%%%%%%%%%%%%%%%%%%%%%%%%%
% Question 5
%%%%%%%%%%%%%%%%%%%%%%%%%%%%%%%%%%%%%%%%%%%%%%%%%%%%%%%%%%%%%%%%%%%%%%%%%%%%%%%%%%%%%%%%
\section*{(Directed) Problem 5: Feedback}
The directed portion took around 10 hours. Mostly, it was getting setup and running scripts, which was done asynchronously for me while doing other stuff. \\ \\


%%TODO
The open ended portion took 3 hours. It was a challenge to develop a program that executes a lot of branch instructions while thinking about how the compiler would generate assembly code. \\ \\
This lab was fun, but I wish there were more coding to do. The directed portion felt like a data analysis section, but it was actually fun analyzing the different benchmarks. The open ended was by far the most fun, because you get exposed to the whole stack from high level C code to low level assembly. I especially had fun working with the RISC-V gnu compiler and seeing how my C code translates to RISC-V, as I've been used to analyzing C code translating to x86. 
I learnt that there are many ways evaluating CPU performance as well as observing and justifying trade-off between certain design decisions. It's very interesting to see the close relationship and connections you can make between an architecture and software, as a developer. \\ \\
There's nothing in this lab that needs to be changed, I think it was a good first lab. \\ \\
Contribution was done individually for each directed part. Then we shared work evenly for open ended, each trying different approaches. I went to OH and found some of the input from other classmates and Josh helpful. 
My partner Kevin Chang helped me with knowledge about how compilers worked and gave me ideas on what kind of algorithms or C syntax can exploit branch type instructions. Every member of the team contributed fairly.

\end{document}